\renewcommand{\prevlecture}{8 }
\renewcommand{\thislecture}{9 }
\renewcommand{\nextlecture}{10 }

%
% Cover page
%

\title[PHYS 201 / Lecture \thislecture]
{
  PHYS 201 / Lecture \thislecture\\
  {\it Electromagnetic Waves}\\
}

\author[C.Andreopoulos] {
  Professor Costas Andreopoulos\inst{1,2}, {\it FHEA}
}
\institute[Liverpool/STFC-RAL] {
   \inst{1} University of Liverpool, Department of Physics\\
   \vspace{0.1cm}
   \inst{2} U.K. Research \& Innovation (UKRI), Science \& Technology Facilities Council,\\
            Rutherford Appleton Laboratory, Particle Physics Department\\
   \vspace{0.5cm}
   {\it {\color{magenta} Lectures delivered at the University of Liverpool, 2021-22}}\\
   \vspace{0.2cm}
}
\date{\today}

\titlegraphic{
  \includegraphics[height=25px]{./images/logo/liverpool.png}
  \hspace{3px}
  \includegraphics[height=30px]{./images/logo/ral.png}
}


\begin{frame}[plain]
  \titlepage
\end{frame}


% ------------------------------------------------------------------------------
% ------------------------------------------------------------------------------

%
% Revision of previous lecture
%

\renewcommand{\lecturesummarytitle}{Revision }

\renewcommand{\summarizedlecture}{8 }


%
%
%

\begin{frame}{Lecture \summarizedlecture revision (Conductors moving in a magnetic field)}

\begin{columns}
  \begin{column}{0.50\textwidth}
     We consider a conductor with length L moves with velocity $\vec{u}$ inside a
     homogenous magnetic field $\vec{B}$, as shown on the right.
  \end{column}
  \begin{column}{0.50\textwidth}
   \begin{center}
     \includegraphics[width=0.90\textwidth]{./images/schematics/conductor_in_magnetic_field_induced_emf.png}\\
   \end{center}
  \end{column}
\end{columns}

\vspace{0.2cm}

Each electron in the conductor feels a magnetic force $\vec{F}_{M} = q \vec{u} \times \vec{B}$.\\
\vspace{0.1cm}
That magnetic force {\bf induces the build-up of charge}
which {\bf produces an electric field} $\vec{E}$:
Each electron feels an electric force $\vec{F}_{E} = q \vec{E}$.\\
\vspace{0.2cm}
The resulting {\bf electric force} $\vec{F}_{E}$ {\bf opposes the magnetic force} $\vec{F}_{M}$.\\

\vspace{0.2cm}

An electrical potential difference develops between the ends of the moving conductor,
which becomes a source of EMF:
\begin{equation*}
   \mathcal{E} = \int_{L} \vec{E} \cdot d\vec{\ell} = \int_{L} \Big( \vec{u} \times \vec{B} \Big) \cdot d\vec{\ell}
\end{equation*}

\end{frame}

%
%
%

\begin{frame}{Lecture \summarizedlecture revision  (Circuit moving in a magnetic field)}

A rectangular circuit moving through a region with magnetic field $\vec{B}$:
\begin{center}
  \includegraphics[width=0.88\textwidth]{./images/schematics/circuit_moving_through_magnetic_field_all_3.png}
\end{center}

\begin{columns}
  \begin{column}{0.55\textwidth}
  In summary:
  {\small
    \setlength{\extrarowheight}{10pt}
    \setlength{\arraycolsep}{5pt}
    \begin{table}[H]
        \begin{tabular}{|c||c|c|}
        \hline
               & $\displaystyle \oint_{L} \vec{E} \cdot d\vec{\ell}$ & $\displaystyle \frac{d\Phi_{M}}{dt}$\\
        \hline
          left   &  uBb &  -uBb \\
          centre & 0    &  0   \\
          right  &  -uBb & uBb \\
        \hline
        \end{tabular}
    \end{table}
  }
  \end{column}
  \begin{column}{0.45\textwidth}
     So, indeed, in all cases:
     \begin{equation*}
       \mathcal{E} = \oint_{L} \vec{E} \cdot d\vec{\ell} = - \frac{d\Phi_{M}}{dt}
     \end{equation*}
  \end{column}
\end{columns}

\end{frame}

%
%
%

\begin{frame}{Lecture \summarizedlecture revision (Faraday's observations)}

In 1831 Michael Faraday reported on a series of experiments.\\
\vspace{0.2cm}

A current flows in a wire loop when:
\vspace{0.1cm}
\begin{enumerate}[(a)]
{\small
  \item the loop is pulled through a magnetic field,
  \item the loop is at rest but the magnet moves in the opposite direction, and
  \item both the loop and the magnet are at rest but the strength of the magnetic field is varied
}
\end{enumerate}

\begin{center}
  \includegraphics[width=0.80\textwidth]{./images/schematics/faraday_law_schematic_with_notes.png}
\end{center}
This led Faraday to realize that:
\begin{center}
{\bf A time-varying magnetic field induces an electic field}
\end{center}

\end{frame}

%
%
%

\begin{frame}{Lecture \summarizedlecture revision (Faraday's law / Lenz's law)}

In all cases the {\bf motional EMF} is directly related
to the {\bf change of the magnetic flux $\Phi_{M}$ though the circuit}:\\

\vspace{0.1cm}

\begin{columns} [T]
  \begin{column}{0.25\textwidth}
    \includegraphics[width=0.98\textwidth]{./images/people/faraday.png}\\
    {\tiny Michael Faraday (1791 - 1867).}
  \end{column}
  \begin{column}{0.75\textwidth}
   \begin{center}
    \begin{equation*}
      \mathcal{E} = \oint_{L} \vec{E} \cdot d\vec{\ell} = - \frac{d\Phi_{M}}{dt}
    \end{equation*}
    \vspace{0.2cm}
    This is the so-called {\bf Faraday's law}.
   \end{center}
  \end{column}
\end{columns}

\vspace{0.1cm}

\begin{columns}
  \begin{column}{0.75\textwidth}
  {\scriptsize
    Lenz' law (1845): The EMF induced by a changing flux has a polarity
    such that the current flowing gives rise to a flux which opposes the change of flux.\\
    \vspace{0.1cm}
    The minus sign is a consequence of the {\bf conservation of energy} and
    of {\bf Newton's 3rd law} of motion: Induction is a an ``inertial reaction''.
    The system develops a current which tries to maintain the flux constant.\\
  }
  \end{column}
  \begin{column}{0.25\textwidth}
    \includegraphics[width=0.83\textwidth]{./images/people/lenz.jpg}\\
    {\tiny Heinrich Lenz (1804 - 1865).}
  \end{column}
\end{columns}

\end{frame}

%
%
%

\begin{frame}{Lecture \summarizedlecture revision (Maxwell correction in Ampere's law)}

We also studied a case where Ampere's law led to paradoxical results.\\

\vspace{0.2cm}

We also saw that Ampere's law (as we knew it) was inconsistent with
the continuity equation (which expresses the local conservation of charge).\\

\vspace{0.2cm}

The problem of course was that we took a law from magnetostatics and
applied it in a different context (electrodynamics) where it is no longer valid.

\vspace{0.2cm}

Maxwell realised that all it takes to fix Ampere's law is to do the following substitution:
\begin{equation*}
  \vec{j} \rightarrow
  \vec{j} + \epsilon_0 \frac{\partial \vec{E}}{\partial t}
\end{equation*}

The term $\displaystyle \epsilon_0 \frac{\partial \vec{E}}{\partial t}$ is the
density of the so-called {\bf displacement current}.

\end{frame}

%
%
%

\begin{frame}{Lecture \summarizedlecture revision (Maxwell's eqs for the dynamic case)}

Compared with what we had seen in the study of electrostatics and magnetostatics,
the study of time-dependent fields (electrodynamics) brought the following complication:\\

\vspace{0.1cm}

\begin{itemize}
   \item {\bf Electric fields are produced} not only by electric charges,
             but also {\bf by changing magnetic fields!}
     \begin{equation*}
        \vec{\nabla} \times \vec{E} = {\color{red} - \frac{\partial \vec{B}}{\partial t}}
     \end{equation*}

   \item {\bf Magnetic fields are produced} not only by electric currents,
             but also {\bf by changing electric fields!}
     \begin{equation*}
         \vec{\nabla} \times \vec{B} = \mu_{0} \Big( \vec{j} + {\color{red} \epsilon_0 \frac{\partial \vec{E}}{\partial t} } \Big)
     \end{equation*}
\end{itemize}

{\small
The full list of Maxwell equations for the static and dynamic cases in vacuum is shown on the next slide.
Notice that all 4 equations are coupled in the dynamic case.\\
}

\end{frame}

%
%
%

\begin{frame}{Lecture \summarizedlecture revision (Maxwell's eqs. for the dynamic case)}

{\small

\begin{center}
{
  \begin{table}[H]
    \begin{tabular}{|l|c|c|}
      \hline
        \multicolumn{3}{|l|} {
          {\color{magenta}
           {\bf Static case (in vacuum)}
          }
        }\\
      \hline
      {\bf Gauss's law} &
        $\displaystyle \oint \vec{E} \cdot d\vec{S} = \frac{1}{\epsilon_0} \int \rho d\tau$ &
        $\displaystyle \vec{\nabla} \cdot \vec{E} = \frac{\rho}{\epsilon_0}$ \\

      {\bf Circuital law} &
        $\displaystyle \oint \vec{E} \cdot d\vec{\ell} = 0$ &
        $\displaystyle \vec{\nabla} \times \vec{E} = 0$ \\

      {\bf Gauss's law} (magn) &
        $\displaystyle  \oint \vec{B} \cdot d\vec{S} = 0$ &
        $\displaystyle  \vec{\nabla} \cdot \vec{B} = 0$ \\

      {\bf Ampere's law} &
        $\displaystyle \oint \vec{B} \cdot d\vec{\ell} = \mu_{0} \int \vec{j} \cdot d\vec{S}$ &
        $\displaystyle \vec{\nabla} \times \vec{B} = \mu_{0} \vec{j}$ \\
      \hline
    \end{tabular}
  \end{table}
}
\end{center}


\begin{center}
{
  \begin{table}[H]
    \begin{tabular}{|l|c|c|}
      \hline
        \multicolumn{3}{|l|} {
          {\color{magenta}
           {\bf Generalization of above for the dynamic case (in vacuum)}
          }
        }\\
      \hline
      {\bf Gauss's law} &
        $\displaystyle \oint \vec{E} \cdot d\vec{S} = \frac{1}{\epsilon_0} \int \rho d\tau$ &
        $\displaystyle \vec{\nabla} \cdot \vec{E} = \frac{\rho}{\epsilon_0}$ \\

      {\bf Faraday's law} &
        $\displaystyle \oint \vec{E} \cdot d\vec{\ell} =  -\frac{\partial}{\partial t} \int \vec{B} \cdot d\vec{S}$ &
        $\displaystyle \vec{\nabla} \times \vec{E} = -  \frac{\partial \vec{B}}{\partial t}$ \\

      {\bf Gauss's law} (magn) &
        $\displaystyle  \oint \vec{B} \cdot d\vec{S} = 0$ &
        $\displaystyle  \vec{\nabla} \cdot \vec{B} = 0$ \\

      {\bf Ampere's law} &
        $\displaystyle \oint \vec{B} \cdot d\vec{\ell} =
           \mu_{0} \int \Big( \vec{j} + \epsilon_0 \frac{\partial \vec{E}}{\partial t}\Big) \cdot d\vec{S}$ &
        $\displaystyle \vec{\nabla} \times \vec{B} =
           \mu_{0} \Big( \vec{j} + \epsilon_0 \frac{\partial \vec{E}}{\partial t}\Big)$ \\
      \hline
    \end{tabular}
  \end{table}
}
\end{center}

}
\end{frame}



%
%
%

\begin{frame}{Lecture \summarizedlecture - \lecturesummarytitle}

Compared with what we had seen in the study of electrostatics and magnetostatics,
the study of time-dependent fields (electrodynamics) brought the following complication:\\

\vspace{0.1cm}

\begin{itemize}
   \item {\bf Electric fields are produced} not only by electric charges,
             but also {\bf by changing magnetic fields!}
     \begin{equation*}
        \vec{\nabla} \times \vec{E} = {\color{red} - \frac{\partial \vec{B}}{\partial t}}
     \end{equation*}

   \item {\bf Magnetic fields are produced} not only by electric currents,
             but also {\bf by changing electric fields!}
     \begin{equation*}
         \vec{\nabla} \times \vec{B} = \mu_{0} \Big( \vec{j} + {\color{red} \epsilon_0 \frac{\partial \vec{E}}{\partial t} } \Big)
     \end{equation*}
\end{itemize}

{\small
The full list of Maxwell equations for the static and dynamic cases in vacuum is shown on the next slide.
Notice that whereas the 2 equations involving $\vec{E}$ and the 2 equations involving $\vec{B}$ were decoupled in the static case,
all 4 equations are coupled in the dynamic case.\\
}

\end{frame}


% %
% %
% %
%
% \begin{frame}{Maxwell's equations for the static and dynamic cases}
%
% {\small
%
% \begin{center}
% {
%   \begin{table}[H]
%     \begin{tabular}{|l|c|c|}
%       \hline
%         \multicolumn{3}{|l|} {
%           {\color{magenta}
%            {\bf Static case (in vacuum)}
%           }
%         }\\
%       \hline
%       {\bf Gauss's law} &
%         $\displaystyle \oint \vec{E} d\vec{S} = \frac{1}{\epsilon_0} \int \rho d\tau$ &
%         $\displaystyle \vec{\nabla} \cdot \vec{E} = \frac{\rho}{\epsilon_0}$ \\
%
%       {\bf Circuital law} &
%         $\displaystyle \oint \vec{E} d\vec{\ell} = 0$ &
%         $\displaystyle \vec{\nabla} \times \vec{E} = 0$ \\
%
%       no magn. monopoles' &
%         $\displaystyle  \oint \vec{B} d\vec{S} = 0$ &
%         $\displaystyle  \vec{\nabla} \cdot \vec{B} = 0$ \\
%
%       {\bf Ampere's law} &
%         $\displaystyle \oint \vec{B} d\vec{\ell} = \mu_{0} \int \vec{j} d\vec{S}$ &
%         $\displaystyle \vec{\nabla} \times \vec{B} = \mu_{0} \vec{j}$ \\
%       \hline
%     \end{tabular}
%   \end{table}
% }
% \end{center}
%
%
% \begin{center}
% {
%   \begin{table}[H]
%     \begin{tabular}{|l|c|c|}
%       \hline
%         \multicolumn{3}{|l|} {
%           {\color{magenta}
%            {\bf Generalization of above for the dynamic case (in vacuum)}
%           }
%         }\\
%       \hline
%       {\bf Gauss's law} &
%         $\displaystyle \oint \vec{E} d\vec{S} = \frac{1}{\epsilon_0} \int \rho d\tau$ &
%         $\displaystyle \vec{\nabla} \cdot \vec{E} = \frac{\rho}{\epsilon_0}$ \\
%
%       {\bf Circuital law} &
%         $\displaystyle \oint \vec{E} d\vec{\ell} =  -\frac{\partial}{\partial t} \int \vec{B} d\vec{S}$ &
%         $\displaystyle \vec{\nabla} \times \vec{E} = -  \frac{\partial \vec{B}}{\partial t}$ \\
%
%       no magn. monopoles &
%         $\displaystyle  \oint \vec{B} d\vec{S} = 0$ &
%         $\displaystyle  \vec{\nabla} \cdot \vec{B} = 0$ \\
%
%       {\bf Ampere's law} &
%         $\displaystyle \oint \vec{B} d\vec{\ell} = \mu_{0} \int \Big( \vec{j} + \epsilon_0 \frac{\partial \vec{E}}{\partial t}\Big) d\vec{S}$ &
%         $\displaystyle \vec{\nabla} \times \vec{B} = \mu_{0} \Big( \vec{j} + \epsilon_0 \frac{\partial \vec{E}}{\partial t}\Big)$ \\
%       \hline
%     \end{tabular}
%   \end{table}
% }
% \end{center}
%
% }
% \end{frame}


%
% Plan for this lecture
%

\begin{frame}{Plan for Lecture \thislecture}

\begin{itemize}
   \item Maxwell's equations in vacuum in absence of sources
   \item Reminder on waves
   \item The electric and magnetic fields satisfy the wave equation
   \item Electromagnetic waves and light
   \item Energy carried by electromagnetic waves
   \item The Poynting theorem
\end{itemize}

\end{frame}

% ------------------------------------------------------------------------------
% ------------------------------------------------------------------------------


%
%
%

\begin{frame}{Maxwell's equations in vacuum in absence of sources}

\begin{columns}
  \begin{column}{0.40\textwidth}
   {\small
       Let me start from the known Maxwell's equations in vacuum,
       and consider the case where there are:
       \begin{itemize}
       {\small
             \item {\bf no charges} ($\rho = 0$)
             \item and {\bf no currents} ($\vec{j} = \vec{0}$).
       }
       \end{itemize}
  }
  \end{column}
  \begin{column}{0.60\textwidth}
     \begin{equation*}
       \vec{\nabla} \vec{E} = \cancelto{0}{\frac{\rho}{\epsilon_0}} \Rightarrow
       {\color{magenta} \vec{\nabla} \vec{E} = 0 }
     \end{equation*}
     \begin{equation*}
        {\color{magenta} \vec{\nabla} \times \vec{E} = - \frac{\partial \vec{B}}{\partial t} }
     \end{equation*}
     \begin{equation*}
        {\color{magenta} \vec{\nabla} \vec{B} = 0 }
     \end{equation*}
     \begin{equation*}
       \vec{\nabla} \times \vec{B} = \cancelto{0}{\mu_0 j} + \epsilon_0 \mu_0 \frac{\partial \vec{E}}{\partial t} \Rightarrow
        {\color{magenta} \vec{\nabla} \times \vec{B} = \epsilon_0 \mu_0 \frac{\partial \vec{E}}{\partial t} }
     \end{equation*}
  \end{column}
\end{columns}

\vspace{0.5cm}

{\bf A change in one field feeds the other} (even in absence of sources):\\
Faraday's law tells us that a changing magnetic field generates an electric field, and
Maxwell's correction to Ampere's law tells us that a changing electric field generates a magnetic field.\\

\vspace{0.2cm}

{\bf Can this interplay generate waves?}

\end{frame}



% starting reminder
{
\reminderslide

%
%
%

\begin{frame}{Reminder: Waves}

\begin{center}
 A wave is a {\bf disturbance that travels through space.}
\end{center}

\begin{columns}
  \begin{column}{0.70\textwidth}
    \includegraphics[width=0.98\textwidth]{./images/photos/mech_wave_water_1.jpg}\\
  \end{column}
  \begin{column}{0.30\textwidth}
    \begin{center}
       As this disturbance travels it transfers energy.\\
       \vspace{0.4cm}
       If there is no absorption or dispersion, that disturbance
       propagates with a constant velocity and a fixed `shape'.
    \end{center}
  \end{column}
\end{columns}

\end{frame}

%
%
%

\begin{frame}{Reminder: Types of waves}

There are {\bf three main types of waves}:

\begin{itemize}

\item {\bf Mechanical waves}\\
    \begin{itemize}
        \item These are the waves you are most familiar with
           \begin{itemize}
                \item For example, sound waves, sea waves, seismic waves.
           \end{itemize}
        \item Mechanical waves {\bf need a medium} to propagate into. They propagate by deforming that medium.
           \begin{itemize}
                 \item When I speak, I disturb the air molecules in front of my mouth.
                 \item They collide with neighbouring molecules and bounce back.
                 \item This disturbance propagates, and it reaches and vibrates your eardrums.
           \end{itemize}
%         \item For mechanical waves, no medium = no propagation.
    \end{itemize}

\item {\bf Electromagnetic waves} {\color{red} $\leftarrow$ The subject of this lecture}\\
    \begin{itemize}
        \item Less familiar, probably, although you interact with them constantly.
        \item They {\bf do not require a medium to propagate.}
    \end{itemize}

\item {\bf Matter waves}  {\color{red} $\leftarrow$ Will study in future modules}\\
    \begin{itemize}
          \item You are probably very unfamiliar with these, but they take centre stage
                    in quantum mechanics ({\em wave-particle duality}).
          \item Waves associated with fundamental particles (such as the electron),
                    composite particles (such as the proton), atoms or even molecules.
    \end{itemize}

\end{itemize}

\end{frame}

%
%
%
%

\begin{frame}{Reminder: The wave equation}

A {\bf wave equation} describes how the disturbance propagates in time.\\

\vspace{0.3cm}

Let $\phi(\vec{r}, t)$ be a function that describes that disturbance as a function of position in space and time.
Then it satisfies the following equation:
\begin{equation*}
   \vec{\nabla}^{2} \phi(\vec{r}, t) = \frac{1}{u^2} \frac{\partial^{2} \phi(\vec{r}, t)} {\partial t^{2}}
\end{equation*}
where {\bf u is the wave velocity}.\\

%\vspace{0.3cm}
%This equation can be derived in many different physical situations and I am sure that you have seen it already in other modules.\\
\vspace{0.3cm}

A mathematician would say that the wave equation is a
{\bf 2$^{nd}$ order linear partial differential equation (p.d.e)}.
\begin{itemize}
   \item This particular type of p.d.e. is called {\em hyperbolic}.
\end{itemize}

\end{frame}

%
%
%
%

\begin{frame}{Reminder: Waves}

Look carefully at the wave equation:
\begin{equation*}
   \vec{\nabla}^{2} \phi(\vec{r}, t) = \frac{1}{u^2} \frac{\partial^{2} \phi(\vec{r}, t)} {\partial t^{2}}
\end{equation*}

Here, we encounter $\vec{\nabla}^{2}$ again.
This is the so-called {\em Laplace operator} that we have already seen at a previous lecture (Poisson equation).\\

\vspace{0.2cm}

The Laplacian $\vec{\nabla}^{2} \phi$ is defined as
{\bf the divergence of the gradient} of the scalar function $\phi$:
\begin{equation*}
   \vec{\nabla}^{2} \phi(\vec{r}, t) =
   \vec{\nabla} \cdot \Big( \vec{\nabla} \phi(\vec{r}, t) \Big) =
   \frac{\partial^{2} \phi(\vec{r}, t)} {\partial x^{2}} +
   \frac{\partial^{2} \phi(\vec{r}, t)} {\partial y^{2}} +
   \frac{\partial^{2} \phi(\vec{r}, t)} {\partial z^{2}}
\end{equation*}

Recall that it represents a quantity which is important in several physical processes:
The Laplacian $\vec{\nabla}^{2} \phi(\vec{r})$ of a scalar function $\phi$ at a point $\vec{r}$
tells you how much $\phi(\vec{r})$ differs from its average over a small volume around $\vec{r}$.

\end{frame}


%
%
%
%

\begin{frame}{Reminder: Solutions of the wave equation}

In one dimension, say x, the wave equation becomes:
\begin{equation*}
    \frac{\partial^{2} \phi(x, t)} {\partial x^{2}} =
    \frac{1}{u^2} \frac{\partial^{2} \phi(x, t)} {\partial t^{2}}
\end{equation*}
where u is the wave velocity.\\

\vspace{0.2cm}

{\bf What are the solutions of the wave equation?}

\vspace{0.2cm}

\underline{\bf Any} (any whatsoever) function $\phi$ of x and t,
{\bf where x and t appear only in the combination of x-ut},
is a solution.
\begin{equation*}
  \phi(x,t) = \psi(z) = \psi(x-ut)
\end{equation*}

\vspace{0.1cm}

For example, the following function is a solution of the wave equation:
\begin{equation*}
  \phi = sin^3\Big( x-ut \Big) - \frac{e^{(x-ut)}}{cos^{\frac{7}{2}}\Big( x-ut\Big)} + \Big( x-ut\Big)^{1.298}
\end{equation*}

\vspace{0.1cm}

We can easily see that using the chain rule.

\end{frame}

%
%
%
%

\begin{frame}{Reminder: Solutions of the wave equation}

Any function of x and t, where x and t appear only in the combination of x-ut,
is a solution of the wave equation.\\

\vspace{0.1cm}

Let $\psi(z)$ be a function of x-ut (z = x-ut).
The partial derivative of  $\psi(z)$ with respect to time is:

\begin{equation*}
  \frac{\partial \psi(z)}{\partial t} =
  \frac{\partial z}{\partial t} \cdot \frac{\partial \psi(z)}{\partial z} =
  \bigg\{ \frac{\partial}{\partial t} \Big( x - ut \Big) \bigg\} \cdot \frac{\partial \psi(z)}{\partial z} =
  -u \frac{\partial \psi(z)}{\partial z}
\end{equation*}

Differentiating once again, we have:
\begin{equation*}
  \frac{\partial^{2} \psi(z)}{\partial t^{2}} =
  \frac{\partial}{\partial t} \Big( \frac{\partial \psi(z)}{\partial t} \Big) =
  \frac{\partial}{\partial t} \Big(  -u \frac{\partial \psi(z)}{\partial z} \Big) =
  \frac{\partial z}{\partial t} \frac{\partial}{\partial z} \Big(  -u \frac{\partial \psi(z)}{\partial z} \Big) =
\end{equation*}

\begin{equation*}
 = \bigg\{ \frac{\partial}{\partial t} \Big( x - ut \Big) \bigg\}  \frac{\partial}{\partial z} \Big(  -u \frac{\partial \psi(z)}{\partial z} \Big) =
  -u \frac{\partial}{\partial z} \Big(  -u \frac{\partial \psi(z)}{\partial z} \Big) =
   u^2 \frac{\partial^{2} \psi(z)}{\partial z^{2}} \Rightarrow
\end{equation*}

{\color{magenta}
\begin{equation*}
  \frac{1}{u^2} \frac{\partial^{2} \psi(z)}{\partial t^{2}} = \frac{\partial^{2} \psi(z)}{\partial z^{2}} (*)
\end{equation*}
}

\end{frame}

%
%
%
%

\begin{frame}{Reminder: Solutions of the wave equation}

Similarly, the partial derivative of  $\psi(z)$ with respect to x is:

\begin{equation*}
  \frac{\partial \psi(z)}{\partial x} =
  \frac{\partial z}{\partial x} \cdot \frac{\partial \psi(z)}{\partial z} =
  \bigg\{ \frac{\partial}{\partial x} \Big( x - ut \Big) \bigg\} \cdot \frac{\partial \psi(z)}{\partial z} =
  \frac{\partial \psi(z)}{\partial z}
\end{equation*}

Differentiating once again, we have:
\begin{equation*}
  \frac{\partial^{2} \psi(z)}{\partial x^{2}} =
  \frac{\partial}{\partial x} \Big( \frac{\partial \psi(z)}{\partial x} \Big) =
  \frac{\partial}{\partial x} \Big( \frac{\partial \psi(z)}{\partial z} \Big) =
  \frac{\partial z}{\partial x} \frac{\partial}{\partial z} \Big( \frac{\partial \psi(z)}{\partial z} \Big) =
\end{equation*}
\begin{equation*}
  =  \bigg\{ \frac{\partial}{\partial x} \Big( x - ut \Big) \bigg\} \frac{\partial}{\partial z} \Big( \frac{\partial \psi(z)}{\partial z} \Big) =
   \frac{\partial^{2} \psi(z)}{\partial z^{2}} \Rightarrow
   {\color{magenta} \frac{\partial^{2} \psi(z)}{\partial x^{2}} =  \frac{\partial^{2} \psi(z)}{\partial z^{2}} } (**)
\end{equation*}

The right-hand sides of Eqs. (*) and (**) are equal, and so are the left-hand sides:
\begin{equation*}
   {\color{blue} \frac{\partial^{2} \psi(z)}{\partial x^{2}} =   \frac{1}{u^2} \frac{\partial^{2} \psi(z)}{\partial t^{2}} }
\end{equation*}

Therefore,  $\psi(z) =  \psi(x-ut)$ is a solution of the wave equation.

\end{frame}


%
%
%
%

\begin{frame}{Reminder: Solutions of the wave equation}

So, any function $\phi$ of x and t,
{\bf where x and t appear only in the combination of x-ut},
is a solution.
\begin{equation*}
  \phi(x,t) = \psi(x-ut)
\end{equation*}

\vspace{0.1cm}

Can you think of {\bf another family of solutions?}\\

\vspace{0.1cm}

Any function $\phi$ of x and t,
{\bf where x and t appear only in the combination of x+ut},
is also a solution.
\begin{equation*}
  \phi(x,t) = \chi(x+ut)
\end{equation*}

\vspace{0.1cm}

The wave equation is {\bf linear}: A superimposition of solutions is a solution.\\

\vspace{0.1cm}

Therefore, most general solution to the wave equation can be written as:\\

\begin{equation*}
  \phi(x,t) = \psi(x-ut) + \chi(x+ut)
\end{equation*}

\end{frame}

%
%
%
%

\begin{frame}{Reminder: Direction of wave propagation}

% See the following two snapshots of a wave taken by a small time interval ${\Delta}t$ apart.
% Imagine that at t=0, the disturbance $\phi(x,0)$ is - Draw?
% At a later fixed time t, the disturbance is $\phi(x,t)$
% But the disturbance is still the same, just displaced in x (because the wave has moved)

What does this particular combination of variables (x-ut, x+ut) tell us about
the direction of wave propagation?
\begin{itemize}
       \item $\psi(x-ut)$ describes a wave moving in the positive direction of x.
       \item $\chi(x+ut)$ describes a wave moving in the negative direction of x.
\end{itemize}

\begin{columns}
  \begin{column}{0.60\textwidth}
     \begin{center}
         \includegraphics[width=0.98\textwidth]{./images/schematics/wave_direction.png}\\
     \end{center}
  \end{column}
  \begin{column}{0.40\textwidth}
  {\scriptsize
     Convince yourselves that the guy on the left should ride a
     $\psi(x-ut)$ wave, not a $\chi(x+ut)$ one.
  }
  \end{column}
\end{columns}

\end{frame}

}% ending reminder


%
%
%
%

\begin{frame}{The electric field satisfies a wave equation}

Now that we have refreshed our memory on waves, let me return to my
original question. Maxwell's equation in vacuum, in the absence of
charges ($\rho = 0$) and currents ($\vec{j} = \vec{0}$) are shown below:

     \begin{equation*}
       \vec{\nabla} \cdot \vec{E} = 0
     \end{equation*}
     \begin{equation*}
       \vec{\nabla} \times \vec{E} = - \frac{\partial \vec{B}}{\partial t}
     \end{equation*}
     \begin{equation*}
       \vec{\nabla} \cdot \vec{B} = 0
      \end{equation*}
      \begin{equation*}
        \vec{\nabla} \times \vec{B} = \epsilon_0 \mu_0 \frac{\partial \vec{E}}{\partial t}
      \end{equation*}

\vspace{0.4cm}

{\bf A change in one field feeds the other} (even in absence of sources).\\

\vspace{0.2cm}

{\bf Can this interplay generate waves?}

\end{frame}

%
%
%
%

\begin{frame}{The electric field satisfies a wave equation}

We will start from Faraday's law in vacuum:
\begin{equation*}
  \vec{\nabla} \times \vec{E} = -\frac{\partial \vec{B}}{\partial t}
\end{equation*}

We will take the curl of both sides and use identity shown below.
Then, we will use Gauss's and Ampere's law in vacuum and in absence of sources:
\begin{equation*}
  \vec{\nabla} \cdot \vec{E} =
       \cancelto{0}{\frac{\rho}{\epsilon_0}} = 0
  \;\;\; and \;\;\;
  \vec{\nabla} \times \vec{B} =
        \cancelto{0}{\mu_0 \vec{j}} + \mu_0 \epsilon_0 \frac{\partial \vec{E}}{\partial t} =
        \mu_0 \epsilon_0 \frac{\partial \vec{E}}{\partial t}
\end{equation*}
to substitute the resulting $\vec{\nabla} \cdot \vec{E}$  and $\vec{\nabla} \times \vec{B}$.

\begin{blockminirem}{\small Mini reminder from calculus}
{\small
For any vector field $\vec{F}$:
$\displaystyle \vec{\nabla} \times \Big( \vec{\nabla} \times \vec{F} \Big) = \vec{\nabla} \Big( \vec{\nabla} \vec{F} \Big) - \vec{\nabla}^{2} \vec{F}$\\
\vspace{0.1cm}
The {\bf curl of the curl} of a vector field is
the {\bf gradient of the divergence} of the field minus
the {\bf divergence of the gradient} (Laplacian) of the field.\\
}
\end{blockminirem}

\end{frame}

%
%
%
%

\begin{frame}{The electric field satisfies a wave equation}

\begin{equation*}
  \vec{\nabla} \times \vec{E} = -\frac{\partial \vec{B}}{\partial t} \Rightarrow
  {\color{red} \vec{\nabla} \times} \Big( \vec{\nabla} \times \vec{E} \Big) =
  {\color{red} \vec{\nabla} \times} \Big( -\frac{\partial \vec{B}}{\partial t} \Big) \Rightarrow
\end{equation*}

\begin{equation*}
  \vec{\nabla} \times (\vec{\nabla} \times \vec{E}) =
  - \vec{\nabla} \times \frac{\partial \vec{B}}{\partial t}
  \xRightarrow{ \vec{\nabla} \times \Big( \vec{\nabla} \times \vec{F} \Big) =
                \vec{\nabla} \Big( \vec{\nabla} \vec{F} \Big) - \vec{\nabla}^{2} \vec{F} }
\end{equation*}

\begin{equation*}
  \vec{\nabla} \Big( \vec{\nabla} \vec{E} \Big) - \vec{\nabla}^{2} \vec{E} =
  - \frac{\partial}{\partial t} (\vec{\nabla} \times \vec{B})
  \xRightarrow{
      \vec{\nabla} \vec{E} = \cancelto{0}{\frac{\rho}{\epsilon_0}} = 0
       \;\; and \;\;
      \vec{\nabla} \times \vec{B} =
           \cancelto{0}{\mu_0 \vec{j}} +\mu_0 \epsilon_0 \frac{\partial \vec{E}}{\partial t} =
           \mu_0 \epsilon_0 \frac{\partial \vec{E}}{\partial t}
  }
\end{equation*}

\begin{equation*}
  - \vec{\nabla}^{2} \vec{E} =
  - \frac{\partial}{\partial t} \Big( \mu_0 \epsilon_0 \frac{\partial \vec{E}}{\partial t} \Big)
  \Rightarrow
  {\color{magenta}
     \vec{\nabla}^{2} \vec{E} =
     \mu_0 \epsilon_0 \frac{\partial^{2} \vec{E}}{\partial t^{2}}
  }
\end{equation*}

\end{frame}


%
%
%
%

\begin{frame}{The magnetic field satisfies a wave equation}

The calculation for $\vec{B}$ proceeds along similar lines.\\
\vspace{0.2cm}

We will start from Ampere's law in vacuum and in absence of sources:
\begin{equation*}
  \vec{\nabla} \times \vec{B} =
        \cancelto{0}{\mu_0 \vec{j}} + \mu_0 \epsilon_0 \frac{\partial \vec{E}}{\partial t} =
        \mu_0 \epsilon_0 \frac{\partial \vec{E}}{\partial t}
\end{equation*}

We will take the curl of both sides and use identity used before.
Then, we will use the fact that there are no magnetic monopoles and Faraday's law:
Gauss's and Ampere's law in vacuum and in absence of sources:
\begin{equation*}
  \vec{\nabla} \cdot \vec{B} = 0
  \;\;\; and \;\;\;
  \vec{\nabla} \times \vec{E} = - \frac{\partial \vec{B}}{\partial t}
\end{equation*}
to substitute the resulting $\vec{\nabla} \cdot \vec{B}$  and $\vec{\nabla} \times \vec{E}$.


\end{frame}


%
%
%
%

\begin{frame}{The magnetic field satisfies a wave equation}

\begin{equation*}
  \vec{\nabla} \times \vec{B} =
  \mu_0 \epsilon_0 \frac{\partial \vec{E}}{\partial t}
  \Rightarrow
  {\color{red} \vec{\nabla} \times} \Big( \vec{\nabla} \times \vec{B} \Big) =
  {\color{red} \vec{\nabla} \times} \Big( \mu_0 \epsilon_0 \frac{\partial \vec{E}}{\partial t} \Big)
  \Rightarrow
\end{equation*}

\begin{equation*}
  \vec{\nabla} \times \Big( \vec{\nabla} \times \vec{B} \Big) =
  \mu_0 \epsilon_0  \vec{\nabla} \times \frac{\partial \vec{E}}{\partial t}
  \xRightarrow{ \vec{\nabla} \times \Big( \vec{\nabla} \times \vec{F} \Big) =
                \vec{\nabla} \Big( \vec{\nabla} \vec{F} \Big) - \vec{\nabla}^{2} \vec{F} }
\end{equation*}

\begin{equation*}
  \vec{\nabla} \Big( \vec{\nabla} \vec{B} \Big) - \vec{\nabla}^{2} \vec{B} =
  \mu_0 \epsilon_0  \frac{\partial}{\partial t} \Big( \vec{\nabla} \times \vec{E} \Big)
  \xRightarrow{ \vec{\nabla} \vec{B} = 0 \;\; and \;\;
                \vec{\nabla} \times \vec{E} = - \frac{\partial \vec{B}}{\partial t} }
\end{equation*}

\begin{equation*}
  - \vec{\nabla}^{2} \vec{B} =
  \mu_0 \epsilon_0  \frac{\partial}{\partial t} \Big(  - \frac{\partial \vec{B}}{\partial t} \Big)
  \Rightarrow
  {\color{magenta}
    \vec{\nabla}^{2} \vec{B} = \mu_0 \epsilon_0 \frac{\partial^{2} \vec{B}}{\partial t^{2}}
  }
\end{equation*}

\end{frame}

%
%
%
%

\begin{frame}{A most peculiar wave!}

\begin{itemize}

   \item The electric and magnetic fields satisfy a wave equation.

   \vspace{0.2cm}

   \item Electric and magnetic field ``disturbances'' can propagate in space
             in absence of charges or currents.

   \vspace{0.2cm}

   \item That disturbance {\bf can not die off}!
             \begin{itemize}
             {\small
               \item A time-varying electric field generates a time-varying magnetic field.
               \item A time-varying magnetic field generates a time-varying electric  field.
             }
             \end{itemize}
             {\bf The change of one field feeds the other!}

   \vspace{0.2cm}

   \item An electric wave can't exist without a magnetic one (and vice versa)\\
             They are part of the same phenomenon: {\bf electromagnetic waves}.

   \vspace{0.2cm}

   \item Unlike mechanical waves, electromagnetic waves {\bf don't need a medium to propagate into}!

\end{itemize}

\end{frame}


%
%
%
%

\begin{frame}{The speed of electromagnetic waves in vacuum}

$\vec{E}$ and $\vec{B}$ satisfy the following wave equations:

\begin{equation*}
  \vec{\nabla}^{2} \vec{E} = \mu_0 \epsilon_0 \frac{\partial^{2} \vec{E}}{\partial t^{2}}
\end{equation*}

\begin{equation*}
  \vec{\nabla}^{2} \vec{B} = \mu_0 \epsilon_0 \frac{\partial^{2} \vec{B}}{\partial t^{2}}
\end{equation*}

Recall that earlier we wrote the wave equation as:
\begin{equation*}
   \vec{\nabla}^{2} \phi(\vec{r}, t) = \frac{1}{u^2} \frac{\partial^{2} \phi(\vec{r}, t)} {\partial t^{2}}
\end{equation*}
where u is the wave velocity.\\

\vspace{0.2cm}

Therefore the electric and magnetic waves propagate in vacuum with the
same speed which, from above, can be identified with:
\begin{equation*}
  u = \frac{1}{\sqrt{\mu_0 \epsilon_0}}
\end{equation*}

\end{frame}

%
%
%
%

\begin{frame}{The speed of electromagnetic waves in vacuum}

Can you confirm that $\frac{1}{\sqrt{\mu_0 \epsilon_0}}$ has velocity units?\\

\vspace{0.4cm}

In SI:
\begin{itemize}
   \item the permeability $\mu_0$ is given in $\frac{N}{A^{2}}$, while
   \item the permitivity $\epsilon_0$ is given in $\frac{C^{2}}{N \cdot m^{2}}$.
\end{itemize}

\vspace{0.2cm}

The product $\mu_0 \epsilon_0$ has units of:
\begin{equation*}
 \frac{\cancel{N}}{A^{2}} \cdot \frac{C^{2}}{\cancel{N} \cdot m^{2}} =
 \frac{1}{A^{2}} \cdot \frac{C^{2}}{m^{2}} =
 \frac{1}{\Big(\frac{C}{s}\Big)^{2}} \cdot \frac{C^{2}}{m^{2}} =
 \frac{s^{2}}{\cancel{C^{2}}} \cdot \frac{\cancel{C^{2}}}{m^{2}} =
 \frac{s^{2}}{m^{2}}
\end{equation*}

\vspace{0.2cm}

Indeed, $\frac{1}{\sqrt{\mu_0 \epsilon_0}}$ has SI units of $\frac{m}{s}$.


\end{frame}

%
%
%
%

\begin{frame}{The speed of electromagnetic waves in vacuum}

What is the numerical value of $\displaystyle \frac{1}{\sqrt{\mu_0 \epsilon_0}}$?

\begin{equation*}
    u =
      \frac{1}{\sqrt{\mu_0 \epsilon_0}} =
       \frac{1}{\sqrt{ \Big( 8.854 \cdot 10^{-12} \; \frac{N}{A} \Big) \cdot
                              \Big( 4\pi \cdot 10^{-7} \; \frac{C}{N\cdot m} \Big)
                            }
                   } =
       {\bf 299,792,458 \; m/s}
\end{equation*}

\vspace{0.2cm}

\begin{columns}
  \begin{column}{0.50\textwidth}
    \begin{center}
       \includegraphics[width=0.65\textwidth]{./images/misc/space_sign_speed_limit.png}\\
    \end{center}
  \end{column}
  \begin{column}{0.50\textwidth}
    \begin{center}
       Do you recognize this number?
    \end{center}
  \end{column}
\end{columns}

\end{frame}


%
%
%
%

\begin{frame}{Measurements of the speed of light before Maxwell's time}

At the time that Maxwell was developing his theory, the speed of light was already known.
In fact, it was known from the late 17th century.\\

\vspace{0.2cm}

\begin{columns}
  \begin{column}{0.33\textwidth}
    \includegraphics[width=0.90\textwidth]{./images/schematics/illustration_from_1676_article_on_the_speed_of_light.jpg}\\
  \end{column}
  \begin{column}{0.67\textwidth}
      In 1670's, {\bf Olaf Roemer} measured the speed of light by observing the {\bf period of the moons of Jupiter.}\\

      \begin{itemize}
      {\small
        \item The apparent period at Earth is different than true period due to the velocity of the Earth.
        \item The variation of the apparent period was
                  \begin{equation*}
                        {\Delta}T = T \cdot \frac{2u}{c}
                  \end{equation*}
                  where T the true period, u the velocity of the earth and c the speed of light.
        \item This allowed the speed of light to be determined to $\sim$~30\%.\\
      }
      \end{itemize}
  \end{column}
\end{columns}

\end{frame}

%
%
%
%

\begin{frame}{Measurements of the speed of light before Maxwell's time}

And in 1848-49, a few years before Maxwell published his famous treatise,
{\bf Hippolyte Fizeau} measured the speed of light to $\sim$5\%.\\

\vspace{0.2cm}

Light passes between two teeth of the rotating gear and gets reflected from a mirror.
     If the reflected light again passes between two teeth of the rotating gear,
     the speed of light can be calculated from the gear-mirror distance and the speed of rotation.\\

\vspace{0.2cm}

\begin{columns}
  \begin{column}{0.40\textwidth}
  {\small
     Fizeau determined the speed of light between an intense light source and a mirror about 8 km distant.
     The light source was interrupted by a rotating cogwheel with 720 notches
     that could be rotated at a variable speed of up to hundreds of times a second.\\
  }
  \end{column}
  \begin{column}{0.60\textwidth}
   \begin{center}
    \includegraphics[width=0.98\textwidth]{./images/schematics/fizeau_experiment_absolute_speed_01.png}\\
   \end{center}
  \end{column}
\end{columns}

\end{frame}

%
%
%
%

\begin{frame}{Light is an electromagnetic wave!}

Maxwell (with typical British understatement) said that the result
was {\em ``somewhat of a surprise''} !!\\
\vspace{0.1cm}

\begin{columns}
  \begin{column}{0.40\textwidth}
   \begin{center}
    \includegraphics[width=0.98\textwidth]{./images/people/maxwell.jpg}\\
     {\small
       James Clerk Maxwell\\
       (1831 - 1879)
     }
   \end{center}
  \end{column}
  \begin{column}{0.60\textwidth}
  {\small
       {\em ``I made out the equations before I had any suspicion of the nearness between
       the two values of the velocity of propagation of magnetic effects and of light, so I think
       I have reasons to believe that the magnetic and luminiferous media are identical''.}\\
       \vspace{0.3cm}
       {\em `We can scarcely avoid the inference that light consists in the transverse undulations
       of the same medium which is the cause of electric and magnetic phenomena.}
       (Maxwell discussed his ideas in terms of a model in which the vacuum was like an elastic solid.)
  }
  \end{column}
\end{columns}

\end{frame}


%
%
%
%

\begin{frame}{Light is an electromagnetic wave!}

Let's just think how amazing and unexpected that was!\\

\vspace{0.1cm}

{\bf What was $\epsilon_0$} ? \\
It is the constant of proportionality in Coulomb's law that, connects
charge densities and the resulting electric field $\vec{E}$:
\begin{equation*}
    \vec{E} (\vec{r})= \frac{1}{4 \pi {\bf \epsilon_0}}
       \int d\tau^{\prime} \rho(\pvec{r}') \frac{\vec{r}-\pvec{r}'}{|\vec{r}-\pvec{r}'|^3}
\end{equation*}

\vspace{0.05cm}

{\bf What was $\mu_0$} ? \\
It is the constant of proportionality in the Biot-Savart law that
connects currents and the resulting magnetic field $\vec{B}$:
\begin{equation*}
   \vec{B} (\vec{r}) = \frac{{\bf \mu_0}}{4\pi}
       \int d\tau^{\prime} \vec{j}(\pvec{r}') \times \frac{\vec{r}-\pvec{r}'}{|\vec{r}-\pvec{r}'|^3}
\end{equation*}

The fact that the speed of the speed of light is determined by these two (apparently
unrelated to light) constants must have stunned Maxwell (or, in his
works, it was {\em ``somewhat of a surprise''}).

\end{frame}


%
%
%
%


\begin{frame}{The electromagnetic spectrum}

\begin{center}
\includegraphics[width=0.98\textwidth]{./images/schematics/emspectrum.png}\\
\end{center}

\end{frame}



%
% Worked example
%

{
\problemslide

%
%
%
%

\begin{frame}{Worked example}

\begin{blockexmplque}{Question}
An electric field $\vec{E}$ is given by:
\begin{equation*}
  \vec{E} = E_{0} \cdot sin\Big(k \cdot  (x-ct)\Big) \cdot \hat{y}
\end{equation*}
where k is the propagation number.\\
\vspace{0.3cm}
\begin{enumerate}[(a)]
 \item Show that it fulfils the wave equation.
 \item Find the corresponding magnetic field.
 \item Comment on directions and amplitudes.
\end{enumerate}
\end{blockexmplque}

\end{frame}

%
%
%
%

\begin{frame}{Worked example}

First, we will calculate the Laplacian ($\vec{\nabla} \cdot \vec{E}$)
of the given electric field $\vec{E}$:
\begin{equation*}
  \vec{E} = E_{0} sin\Big(k(x-ct)\Big) \hat{y}
\end{equation*}

We have:
\begin{equation*}
  \vec{\nabla}^{2} \vec{E} =
  \frac{\partial^{2} \vec{E}} {\partial x^{2}} +
  \cancelto{0}{\frac{\partial^{2} \vec{E}} {\partial y^{2}}} +
  \cancelto{0}{\frac{\partial^{2} \vec{E}} {\partial z^{2}}} =
  \frac{\partial^{2}} {\partial x^{2}} \bigg\{ E_{0} sin\Big(k(x-ct)\Big) \hat{y} \bigg\} =
\end{equation*}
\begin{equation*}
     = E_{0} \bigg\{ \frac{\partial^{2}} {\partial x^{2}} sin\Big(k(x-ct)\Big) \bigg\} \hat{y} =
        E_{0} \bigg\{ k \frac{\partial} {\partial x} cos\Big(k(x-ct)\Big) \bigg\} \hat{y} =
\end{equation*}
\begin{equation*}
     = E_{0} \bigg\{ - k^2 sin\Big(k(x-ct)\Big) \bigg\} \hat{y} =
      - k^2 \bigg\{ E_{0} sin\Big(k(x-ct)\Big) \hat{y} \bigg\} \Rightarrow
\end{equation*}
\begin{equation*}
  \vec{\nabla}^{2} \vec{E} = - k^2 \vec{E}
\end{equation*}

\end{frame}

%
%
%
%

\begin{frame}{Worked example}

Then, we will calculate the second-order partial derivative with respect to time
($\partial^2 / \partial t^2$) of the given electric field $\vec{E}$:
\begin{equation*}
  \vec{E} = E_{0} sin\Big(k(x-ct)\Big) \hat{y}
\end{equation*}

\begin{equation*}
  \frac{\partial^{2} \vec{E}} {\partial t^{2}} =
  \frac{\partial^{2}} {\partial t^{2}} \bigg\{ E_{0} sin\Big(k(x-ct)\Big) \hat{y} \bigg\} =
\end{equation*}
\begin{equation*}
  = E_{0} \bigg\{ \frac{\partial^{2}} {\partial t^{2}} sin\Big(k(x-ct)\Big) \bigg\} \hat{y} =
     E_{0} \bigg\{-kc \frac{\partial} {\partial t} cos\Big(k(x-ct)\Big) \bigg\} \hat{y} =
\end{equation*}
\begin{equation*}
   = E_{0} \bigg\{ - k^{2} c^{2} sin\Big(k(x-ct)\Big) \bigg\} \hat{y} =
    - k^{2} c^{2} \bigg\{ E_{0} sin\Big(k(x-ct)\Big) \hat{y} \bigg\} \Rightarrow
\end{equation*}
\begin{equation*}
  \frac{\partial^{2} \vec{E}} {\partial t^{2}} = - k^2 c^2 \vec{E}
\end{equation*}

\end{frame}


%
%
%
%

\begin{frame}{Worked example}

So, we found that:
\begin{equation*}
  \vec{\nabla}^{2} \vec{E} = - k^2 \cdot \vec{E}
\end{equation*}

and:
\begin{equation*}
  \frac{\partial^{2} \vec{E}} {\partial t^{2}} = - k^2 \cdot c^2 \cdot \vec{E}
\end{equation*}

Hence:
\begin{equation*}
  \frac{\partial^{2} \vec{E}} {\partial t^{2}} = c^2 \cdot \vec{\nabla}^{2} \vec{E}
\end{equation*}

The electric field $\displaystyle \vec{E} = E_{0} \cdot sin\Big(k \cdot (x-c \cdot t)\Big) \cdot \hat{y}$
fullfils a wave equation with velocity c.

\end{frame}


%
%
%
%

\begin{frame}{Worked example}

To find the corresponding magnetic field $\vec{B}$, we will use Faraday's law:
\begin{equation*}
  \vec{\nabla} \times \vec{E} = - \frac{\partial \vec{B}}{\partial t}
\end{equation*}

First we will calculate $\vec{\nabla} \times \vec{E}$,
and then integrate the result over time.\\

\vspace{0.2cm}

We have:
\begin{equation*}
  \vec{\nabla} \times \vec{E} =
   \left|
    \begin{array}{ccc}
      \hat{x}                     & \hat{y}                     & \hat{z} \\
      \frac{\partial}{\partial x} & \frac{\partial}{\partial y} & \frac{\partial}{\partial z} \\
       E_{x}                      &  E_{y}                      &  E_{z} \\
    \end{array}
   \right| =
   \hat{x}
   \left|
    \begin{array}{cc}
      \frac{\partial}{\partial y} & \frac{\partial}{\partial z} \\
       E_{y}                      &  E_{z} \\
    \end{array}
   \right| -
   \hat{y}
   \left|
    \begin{array}{cc}
      \frac{\partial}{\partial x} & \frac{\partial}{\partial z} \\
       E_{x}                      &  E_{z} \\
    \end{array}
   \right| +
   \hat{z}
   \left|
    \begin{array}{cc}
      \frac{\partial}{\partial x} & \frac{\partial}{\partial y} \\
       E_{x}                      &  E_{y}                      \\
    \end{array}
   \right| \Rightarrow
\end{equation*}

\begin{equation*}
   \vec{\nabla} \times \vec{E} =
   \hat{x} \cdot \Big( \frac{\partial E_{z}}{\partial y} - \frac{\partial E_{y}}{\partial z} \Big) -
   \hat{y} \cdot \Big( \frac{\partial E_{z}}{\partial x} - \frac{\partial E_{x}}{\partial z} \Big) +
   \hat{z} \cdot \Big( \frac{\partial E_{y}}{\partial x} - \frac{\partial E_{x}}{\partial y} \Big)
\end{equation*}

\end{frame}

%
%
%
%

\begin{frame}{Worked example}

Recall that $\vec{E}$ has only a y component ($E_{x} = E_{z} = 0$),
which has only x dependence ($E_{y} = E_{y}(x)$ with no y or z dependence).
Therefore:

\begin{equation*}
  \vec{\nabla} \times \vec{E} =
   \hat{x} \cdot \Big( \cancelto{0}{\frac{\partial E_{z}}{\partial y}} - \cancelto{0}{\frac{\partial E_{y}}{\partial z}} \Big) -
   \hat{y} \cdot \Big( \cancelto{0}{\frac{\partial E_{z}}{\partial x}} - \cancelto{0}{\frac{\partial E_{x}}{\partial z}} \Big) +
   \hat{z} \cdot \Big( \frac{\partial E_{y}}{\partial x} - \cancelto{0}{\frac{\partial E_{x}}{\partial y}} \Big) \Rightarrow
\end{equation*}

\begin{equation*}
  \vec{\nabla} \times \vec{E} = \hat{z} \cdot\frac{\partial E_{y}}{\partial x}
\end{equation*}

The first-order partial derivative of $E_{y}(x)$ with respect to x is:
\begin{equation*}
  \frac{\partial E_{y}}{\partial x} =
      \frac{\partial}{\partial x} \bigg\{ E_{0} \cdot sin\Big(k \cdot (x-c \cdot t)\Big) \bigg\} =
      E_{0} \cdot k \cdot cos\Big(k \cdot (x-c \cdot t)\Big)
\end{equation*}

Hence:
\begin{equation*}
  \vec{\nabla} \times \vec{E} = \bigg\{  E_{0} \cdot k \cdot cos\Big(k \cdot (x-c \cdot t)\Big) \bigg\} \cdot \hat{z}
\end{equation*}

\end{frame}

%
%
%
%

\begin{frame}{Worked example}

Substituting the previous result in Faraday's law, we have:
\begin{equation*}
  - \frac{\partial \vec{B}}{\partial t} = \bigg\{  E_{0} \cdot k \cdot cos\Big(k \cdot (x-c \cdot t)\Big) \bigg\} \cdot \hat{z}
\end{equation*}

Integrating over time:
\begin{equation*}
  \vec{B} = \bigg\{ - E_{0} \cdot k \cdot \int cos\Big(k \cdot (x-c \cdot t)\Big) dt \bigg\} \cdot \hat{z} =
\end{equation*}
\begin{equation*}
   = \bigg\{ - E_{0} \cdot k \cdot \frac{sin\Big(k \cdot (x-c \cdot t)\Big)}{-k \cdot c} \bigg\} \cdot \hat{z} + \vec{C} =
       \bigg\{ \frac{E_{0}}{c} \cdot sin\Big(k \cdot (x-c \cdot t)\Big) \bigg\} \cdot \hat{z} + \vec{C}
\end{equation*}

In the absence of currents there are no time invariant components of $\vec{B}$ ($\vec{C} = \vec{0}$) hence:
\begin{equation*}
  \vec{B} = \bigg\{ \frac{E_{0}}{c} \cdot sin\Big(k \cdot (x-c \cdot t)\Big) \bigg\} \cdot \hat{z}
\end{equation*}


\end{frame}

%
%
%
%

\begin{frame}{Worked example}

Our electric and magnetic fields are given by:
\begin{equation*}
   \vec{E} = \bigg\{ E_{0} \cdot sin\Big(k \cdot (x-c \cdot t)\Big) \bigg\} \cdot \hat{y}
    \;\;\;\; and \;\;\;\;
   \vec{B} = \bigg\{ \frac{E_{0}}{c} \cdot sin\Big(k \cdot (x-c \cdot t)\Big) \bigg\} \cdot \hat{z}
\end{equation*}

Notice that the {\bf amplitude of the magnetic field is smaller to that of the electric field
by a factor of c}.\\

\vspace{0.2cm}

Also notice that:
\begin{itemize}
   \item the wave propagates on the x-axis,
   \item the electric field oscillates on the y-axis, and
   \item the magnetic field oscillates on the z-axis
\end{itemize}

Therefore the {\bf electric and magnetic fields are}:
\begin{itemize}
    \item {\bf perpendicular to each other}, and
    \item {\bf perpendicular to the direction of motion}.
\end{itemize}

\end{frame}

%
%
%
%

\begin{frame}{Worked example}

\begin{center}
\includegraphics[width=0.80\textwidth]{./images/schematics/emwave_01.jpg}\\
\end{center}

\end{frame}

} % Worked example


%
%
%
%

\begin{frame}{Energy carried by electromagnetic waves}

We know that an e/m wave carries energy.\\
\vspace{0.1cm}

How can we express the {\bf energy transported by an e/m wave?}\\
\vspace{0.2cm}

The energy stored in the electric field is:
\begin{equation*}
  U_{e} = \frac{\epsilon_0}{2} \int_{\tau} d\tau |\vec{E}|^2
\end{equation*}
while the energy stored in the magnetic field is:
\begin{equation*}
  U_{m} = \frac{1}{2\mu_0} \int_{\tau} d\tau |\vec{B}|^2
\end{equation*}
Therefore the energy stored in an electromagnetic field is:
\begin{equation*}
  U_{em} = \frac{1}{2} \int_{\tau} d\tau \Big( \epsilon_0 |\vec{E}|^2 + \frac{1}{\mu_0} |\vec{B}|^2 \Big)
\end{equation*}

\end{frame}

%
%
%
%

\begin{frame}{Energy carried by electromagnetic waves}

Consider fields $\vec{E}$ and $\vec{B}$ within some volume $\tau$.
A charge q moves within that volume with velocity $\vec{u}$ and it is displaced by a short distance $d\vec{\ell}$.\\
\vspace{0.2cm}
The rate of change of the work W done on the charge q by the
e/m field\\ (see Optional Reading) is given by:
\begin{equation*}
  \frac{dW}{dt} =
     - {\color{magenta} \oint_{S} d\vec{S} \frac{1}{\mu_0} \Big(\vec{E} \times \vec{B} \Big) }
     - \frac{\partial}{\partial t} \int_{\tau} d\tau  \Big( \frac{1}{2\mu_0} \vec{B}^2 + \frac{\epsilon_0}{2} \vec{E}^2 \Big)
\end{equation*}

As could be anticipated, the {\bf work done on the charge q} is related with the
{\bf decrease of the energy stored in the e/m field} (2nd term on right).\\
\vspace{0.2cm}

The additional {\color{magenta}``surface'' term} above,
can be interpreted as the {\em energy carried away
from the volume $\tau$, through its surface S, by the e/m field}.\\
\vspace{0.2cm}

The above expression is known as the {\bf Poynting theorem}.


\end{frame}


%
%
%
%

\begin{frame}{The Poynting vector}

We define the {\bf Poynting vector} $\vec{N}$ as:
\begin{equation*}
  \vec{N} = \frac{1}{\mu_0} \Big( \vec{E} \times \vec{B} \Big)
\end{equation*}

\vspace{0.3cm}

It represents the {\bf energy flux density (rate energy transfer per unit area)} and it has units of $W(att)/m^2$.\\

\vspace{0.3cm}

Recall that the cross product $\vec{A} \times \vec{B}$ of 2 vectors $\vec{A}$ and $\vec{B}$
is perpendicular both to $\vec{A}$ and $\vec{B}$. So, in an electromagnetic wave, the energy
flows in a direction which is perpendicular to both the electric and magnetic field.

\end{frame}

%
% Worked example
%

{
\problemslide

%
%
%
%

\begin{frame}{Worked example}

An e/m wave is given by the time-dependent fields:
\begin{equation*}
  \vec{E}(\vec{r},t) = E_0 sin\Big( k(z-ct) \Big) \hat{x}
\end{equation*}
\begin{equation*}
  \vec{B}(\vec{r},t) = \frac{E_0}{c} sin\Big( k(z-ct) \Big) \hat{y}
\end{equation*}

This gives the Poynting vector:
\begin{equation*}
  \vec{N} =
    \frac{1}{\mu_0} \Big( \vec{E} \times \vec{B} \Big) =
    \frac{E_0^2}{\mu_0 c} sin^{2}\Big( k(z-ct) \Big) \Big( \hat{x} \times \hat{y} \Big)
\end{equation*}

One can easily see that:
\begin{equation*}
  \hat{x} \times \hat{y} = \hat{z}
\end{equation*}

Therefore, the Poynting vector is:
\begin{equation*}
  \vec{N} =
    \frac{E_0^2}{\mu_0 c} sin^{2}\Big( k(z-ct) \Big) \hat{z}
\end{equation*}

\end{frame}

%
%
%
%

\begin{frame}{Worked example}

Therefore, the Poynting vector is:
\begin{equation*}
  \vec{N} =
    \frac{E_0^2}{\mu_0 c} sin^{2}\Big( k(z-ct) \Big) \hat{z}
\end{equation*}

The above vector describes the  power transmitted by the e/m wave
per unit area perpendicular to its propagation.

The average power transmitted by the e/m wave, is given by
the {\bf average of $\vec{N}$ over a period T}:
\begin{equation*}
  <\vec{N}> = \frac{1}{T} \int_{0}^{T} \vec{N} dt
\end{equation*}

\begin{equation*}
  <\vec{N}> = \frac{1}{T} \int_{0}^{T} \Big( \frac{E_0^2}{\mu_0 c} sin^{2}\Big( k(z-ct) \Big) \hat{z} \Big) dt
\end{equation*}

\end{frame}

%
%
%
%

\begin{frame}{Worked example}

\begin{equation*}
  <\vec{N}> = \frac{1}{T} \int_{0}^{T} \Big( \frac{E_0^2}{\mu_0 c} sin^{2}\Big( k(z-ct) \Big) \hat{z} \Big) dt
\end{equation*}

\begin{equation*}
  <\vec{N}> =  \hat{z} \frac{1}{T} \frac{E_0^2}{\mu_0 c} \int_{0}^{T} sin^{2}\Big( k(z-ct) \Big) dt
\end{equation*}

\vspace{0.2cm}

The average of $sin^{2}\theta$ over a period is 1/2, therefore:
\begin{equation*}
  <\vec{N}> =  \hat{z} \frac{E_0^2}{2 \mu_0 c}
\end{equation*}

\end{frame}

} % A worked example


%
%
%
%

\begin{frame}{Radiation Pressure}

E/M waves {\bf carry momentum}:
Can {\bf exert pressure} on an object when shining on it!\\
\vspace{0.1cm}

Assume that an object
is illuminated by radiation for a time interval ${\Delta}t$ and that
the radiation is entirely {\bf absorbed} by the object.
If its energy gain is ${\Delta}U$, then the magnitude of its momentum
change ${\Delta}p$ is given by:
\begin{equation*}
  {\Delta}p = \frac{{\Delta}U}{c}
  \;\;\; {\color{red} (total \; absorption)}
\end{equation*}
where c is the speed of light.\\
\vspace{0.1cm}

If, instead of being fully absorbed, the radiation is fully reflected back
by the object along its original path, the momentum change of the
object is twice  the one given above:
\begin{equation*}
  {\Delta}p = \frac{2{\Delta}U}{c}
  \;\;\; {\color{red} (total \; reflection \; along \; original \; path)}
\end{equation*}

In realistic situations the momentum change is between $\frac{{\Delta}U}{c}$
and $\frac{2{\Delta}U}{c}$.

\end{frame}

%
%
%
%

\begin{frame}{Radiation Pressure}

From Newton's second law, a change in momentum is related to a force
\begin{equation*}
   F = \frac{{\Delta}p}{{\Delta}t}
\end{equation*}

If the radiation has intensity I (power per area, or energy per time
per area), then over a time interval ${\Delta}t$, the energy
intercepted by an area A is:
\begin{equation*}
  {\Delta}U = I A {\Delta}t
\end{equation*}
Therefore, the force exerted by radiation on an object is given by:
\begin{equation*}
  F = \frac{IA}{c}
  \;\;\; {\color{red} (total \; absorption)}
\end{equation*}
\begin{equation*}
  F = \frac{2IA}{c}
  \;\;\; {\color{red} (total \; reflection \; along \; original \; path)}
\end{equation*}

If the radiation is partially absorbed and partially reflected, the
exerted force will be between $\frac{IA}{c}$
and $\frac{2IA}{c}$.

\end{frame}

%
%
%
%

\begin{frame}{Radiation Pressure}

The force per unit area on an object due to radiation is the {\bf radiation
pressure} $P_{r}$.\\
\vspace{0.2cm}

From the previous expressions, we obtain:
\begin{equation*}
  P_r = \frac{I}{c}
  \;\;\; {\color{red} (total \; absorption)}
\end{equation*}
\begin{equation*}
  P_r = \frac{2I}{c}
  \;\;\; {\color{red} (total \; reflection \; along \; original \; path)}
\end{equation*}


\end{frame}

%
%
%
%

\begin{frame}{Radiation Pressure}

{\bf Solar sail} a form of spacecraft propulsion!
\begin{itemize}
{\small
  \item Using radiation pressure from sunlight (see tomorrow's
    workshop tasks)
  \item See {\color{blue}https://en.wikipedia.org/wiki/Solar\_sail}
}
\end{itemize}

\vspace{0.2cm}

\begin{columns}
  \begin{column}{0.50\textwidth}
   \begin{center}
         \includegraphics[width=0.90\textwidth]{./images/photos/lightsail.jpg}\\
   \end{center}
  \end{column}
  \begin{column}{0.50\textwidth}
       {\bf LighSail 2}
      \begin{itemize}
      {\small
        \item Area: 32 m$^2$
        \item Launch date: March 2017!
        \item See {\color{blue}http://sail.planetary.org}
      }
      \end{itemize}
  \end{column}
\end{columns}

\vspace{0.4cm}

What if we replace sunlight with powerful directed laser beams?\\
See: ``NASA thinks there's a way to get to Mars in three days''\\
{\color{blue}http://phys.org/news/2016-02-nasa-mars-days.html}\\

\end{frame}


%
%
%

\begin{frame}{Radiation Pressure}

Who cares about Mars!\\
\vspace{0.2cm}

Shoot for {\bf Alpha Centauri} and its recently-discovered planet Proxima b!
\vspace{0.2cm}

See {\bf Breakthrough Starshot}:
{\color{blue}http://breakthroughinitiatives.org/Initiative/3}\\

\end{frame}

% ------------------------------------------------------------------------------
% ------------------------------------------------------------------------------

%
% What to remember
%

\renewcommand{\lecturesummarytitle}{Main points to remember }

\renewcommand{\summarizedlecture}{9 }

%
%
%

\begin{frame}{Lecture \summarizedlecture - \lecturesummarytitle}

\begin{itemize}

\item
We looked at Maxwell's equation in vacuum in absence of sources, and we
saw that {\bf the electric and magnetic fields satisfy a wave equation}:
\begin{equation*}
  \vec{\nabla}^{2} \vec{E} = \mu_0 \epsilon_0 \frac{\partial^{2} \vec{E}}{\partial t^{2}}
  \;\;\;\; and \;\;\;\;
  \vec{\nabla}^{2} \vec{B} = \mu_0 \epsilon_0 \frac{\partial^{2} \vec{B}}{\partial t^{2}}
\end{equation*}

\item
The velocity of an electromagnetic wave is given by:
\begin{equation*}
    \frac{1}{\sqrt{\mu_0 \epsilon_0}} = {\bf 299,792,458 \; m/s}  \;\;\;(the \; speed \; of \; light)
\end{equation*}

\item
This led Maxwell to realize that {\bf light is an electromagnetic wave}.

\end{itemize}

\end{frame}

%
%
%

\begin{frame}{Lecture \summarizedlecture - \lecturesummarytitle (cont'd)}

\begin{itemize}

\item
An electromagnetic wave is a {\bf most peculiar wave}!

   \begin{itemize}
   {\small
     \item An electromagnetic wave  {\bf can not die off}!
             \begin{itemize}
             {\scriptsize
               \item A time-varying electric field generates a time-varying magnetic field.
               \item A time-varying magnetic field generates a time-varying electric  field.
             }
             \end{itemize}
             {\bf The change of one field feeds the other!}

     \vspace{0.1cm}

     \item An electric wave can't exist without a magnetic one (and vice versa)\\
               They are part of the same phenomenon: {\bf electromagnetic waves}.

     \vspace{0.1cm}

     \item The electric and magnetic fields are perpendicular to each
               other an perpendicular to the direction of wave propagation.

     \vspace{0.1cm}

     \item Electromagnetic waves {\bf don't need a medium to propagate}!
   }
   \end{itemize}

\item
The {\bf Poynting vector} $\vec{N}$ as:
\begin{equation*}
  \vec{N} = \frac{1}{\mu_0} \Big( \vec{E} \times \vec{B} \Big)
\end{equation*}
It represents the {\bf energy flux density (rate energy transfer per unit area)} and it has units of $W(att)/m^2$.\\

\end{itemize}

\end{frame}


%
%
%

\begin{frame}{Lecture \summarizedlecture - \lecturesummarytitle (cont'd)}

\begin{itemize}

\item
E/M waves {\bf carry momentum}:
Can {\bf exert pressure} on an object when shining on it!\\

\item
The radiation pressure was given by:
\begin{equation*}
  P_r = \frac{I}{c}
  \;\;\; {\color{red} (total \; absorption)}
\end{equation*}
\begin{equation*}
  P_r = \frac{2I}{c}
  \;\;\; {\color{red} (total \; reflection \; along \; original \; path)}
\end{equation*}
where I is the intensity I (power per area, or energy per time
per area) of the radiation.

\end{itemize}

\end{frame}


%
% Plan for the next lecture
%

\begin{frame}{At the next lecture (Lecture \nextlecture)}

\begin{itemize}
\item Electromagnetic waves in matter
\end{itemize}

\end{frame}

%
% Optional reading
%


\begin{frame}[plain,c]
\begin{center}
{\Huge \bf Optional reading for Lecture \thislecture}
\end{center}
\end{frame}

% ------------------------------------------------------------------------------

%
%
%
%

\begin{frame}{Energy carried by electromagnetic waves}

We will see that the  {\em energy carried away from the volume $\tau$ by an e/m wave
through the surface S} is given by:
\begin{equation*}
     \oint_{S} d\vec{S} \cdot \frac{1}{\mu_0} \Big( \vec{E} \times \vec{B} \Big)
\end{equation*}

\vspace{0.3cm}

Let's assume that I have an electric field $\vec{E}$ and a magnetic field $\vec{B}$ within some volume $\tau$.
A charge q moves within that volume with velocity $\vec{u}$ and it is displaced by a short distance $d\vec{\ell}$.
\begin{itemize}
   \item There is work done by the e/m force.
   \item The energy has to come from somewhere.
       \begin{itemize}
            \item Where from? Well, there is the energy
                      $\displaystyle \frac{1}{2} \int_{\tau} d\tau \Big( \epsilon_0 |\vec{E}|^2 + \frac{1}{\mu_0} |\vec{B}|^2 \Big)$
                      stored in the e/m field.
             \item So, if there is work done on the charge, the energy stored in the e/m field must be decreased.
       \end{itemize}
\end{itemize}

\end{frame}

%
%
%
%

\begin{frame}{Energy carried by electromagnetic waves}

The work done by the Lorentz force is:
\begin{equation*}
  dW = \vec{F} \cdot d\vec{\ell} = q \Big( \vec{E} + \vec{u} \times \vec{B} \Big) \cdot \vec{u} dt
\end{equation*}

We have already seen that magnetic forces do no work:
The $\vec{u} \times \vec{B}$ vector is perpendicular to $\vec{u}$ so the dot product
$\Big( \vec{u} \times \vec{B} \Big) \cdot \vec{u}$ is 0.

So the work done can be written just as
\begin{equation*}
  dW = \vec{F} \cdot d\vec{\ell} = q \vec{E} \cdot \vec{u} dt
\end{equation*}

To move from the discrete to the continuous case, take q to be the amount of charge contained
within a volume $d\tau$ within a region characterised by charge density $\rho$.
Then
\begin{equation*}
  dW = q \vec{E} \cdot \vec{u} dt \rightarrow \Big( \rho d\tau \Big) \vec{E} \cdot \vec{u} dt =
           dt d\tau \vec{E} \cdot \Big( \rho \vec{u} \Big)
\end{equation*}


\end{frame}


%
%
%
%

\begin{frame}{Energy carried by electromagnetic waves}

The work done per unit time, in the infinitesimal volume $d\tau$ is:
\begin{equation*}
  \frac{dW}{dt} =  d\tau \vec{E}  \cdot \Big( \rho \vec{u} \Big)
\end{equation*}

As we have seen, the product $\rho \vec{u}$ is the current density $\vec{j}$. Therefore, the work done per unit time and per unit volume can be written as:
\begin{equation*}
  \frac{dW}{dt} =  d\tau \vec{E}  \cdot \vec{j}
\end{equation*}

The total work (integrated over all space) per unit time is:
\begin{equation*}
  \frac{dW}{dt} =  \int_{\tau} d\tau \vec{E} \cdot \vec{j}
\end{equation*}

To proceed, we need to examine the dot product $\vec{E} \cdot \vec{j}$ of the electric field $\vec{E}$ and the current density $\vec{j}$.

\end{frame}

%
%
%
%

\begin{frame}{Energy carried by electromagnetic waves}

Starting from Ampere's law, we can write $\vec{j}$ in terms of
$\vec{\nabla} \times \vec{B}$ and $\frac{\partial \vec{E}}{\partial t}$:
\begin{equation*}
  \vec{\nabla} \times \vec{B} = \mu_0 \Big( \vec{j} + \epsilon_0 \frac{\partial \vec{E}}{\partial t} \Big) \Rightarrow
  \vec{j} = \frac{1}{\mu_0}  \vec{\nabla} \times \vec{B} - \epsilon_0 \frac{\partial \vec{E}}{\partial t}
\end{equation*}

Therefore:
\begin{equation*}
 \vec{E} \cdot \vec{j} =
   \vec{E} \cdot \Big( \frac{1}{\mu_0}  \vec{\nabla} \times \vec{B} - \epsilon_0 \frac{\partial \vec{E}}{\partial t} \Big) =
   \frac{1}{\mu_0}  \vec{E} \cdot \Big( \vec{\nabla} \times \vec{B} \Big) - \epsilon_0 \vec{E} \cdot \frac{\partial \vec{E}}{\partial t}
\end{equation*}

The cross-product term can be written as:
\begin{equation*}
 \vec{E} \cdot \Big( \vec{\nabla} \times \vec{B} \Big) =
 - \vec{\nabla} \cdot \Big( \vec{E} \times \vec{B} \Big) - \vec{B} \cdot \frac{\partial \vec{B}}{\partial t}
\end{equation*}
Try to prove it at home! (Details on the next slide.)

\end{frame}

%
%
%
%

\begin{frame}{Further details (skip on a first read)}

Starting from $\vec{\nabla} \cdot \Big( \vec{E} \times \vec{B} \Big)$,
we can write using the {\em CAB-BAC} identity:
\begin{equation*}
 \vec{\nabla} \cdot \Big( \vec{E} \times \vec{B} \Big) =
   \vec{B} \cdot \Big( \vec{\nabla} \times \vec{E} \Big) - \vec{E} \cdot \Big( \vec{\nabla} \times \vec{B} \Big)
\end{equation*}

The $\vec{E} \cdot \Big( \vec{\nabla} \times \vec{B} \Big)$ term
is what appears in the $\vec{E} \cdot \vec{j}$ dot product.\\
Solving for it:
\begin{equation*}
 \vec{E} \cdot \Big( \vec{\nabla} \times \vec{B} \Big) =
   - \vec{\nabla} \cdot \Big( \vec{E} \times \vec{B} \Big) + \vec{B} \cdot \Big( \vec{\nabla} \times \vec{E} \Big)
\end{equation*}

From Faraday's law:
\begin{equation*}
  \vec{\nabla} \times \vec{E} = - \frac{\partial \vec{B}}{\partial t}
\end{equation*}

Therefore:
\begin{equation*}
 \vec{E} \cdot \Big( \vec{\nabla} \times \vec{B} \Big) =
  - \vec{\nabla} \cdot \Big( \vec{E} \times \vec{B} \Big) - \vec{B} \cdot \frac{\partial \vec{B}}{\partial t}
\end{equation*}

\end{frame}

%
%
%
%

\begin{frame}{Energy carried by electromagnetic waves}


\begin{equation*}
 \vec{E} \cdot \vec{j} =
    - \frac{1}{\mu_0} \vec{\nabla} \cdot \Big( \vec{E} \times \vec{B} \Big)
    - \frac{1}{\mu_0} \vec{B} \cdot \frac{\partial \vec{B}}{\partial t} - \epsilon_0 \vec{E} \cdot \frac{\partial \vec{E}}{\partial t}
\end{equation*}

For a vector field $\vec{F}$, we can easily see that:
\begin{equation*}
  \vec{F} \cdot \frac{\partial \vec{F}}{\partial t} = \frac{1}{2} \frac{\partial |\vec{F}|^2}{\partial t}
\end{equation*}

and, therefore:
\begin{equation*}
 \vec{E} \cdot \vec{j} =
    - \frac{1}{\mu_0} \vec{\nabla} \cdot \Big( \vec{E} \times \vec{B} \Big)
    - \frac{1}{2\mu_0} \frac{\partial |\vec{B}|^2}{\partial t} - \frac{\epsilon_0}{2} \frac{\partial |\vec{E}|^2}{\partial t} \Rightarrow
\end{equation*}

\begin{equation*}
 \vec{E} \cdot \vec{j} =
    - \frac{1}{\mu_0} \vec{\nabla} \cdot \Big( \vec{E} \times \vec{B} \Big)
    - \frac{\partial}{\partial t} \Big( \frac{1}{2\mu_0} |\vec{B}|^2 + \frac{\epsilon_0}{2} |\vec{E}|^2 \Big)
\end{equation*}

\end{frame}


%
%
%
%

\begin{frame}{Energy carried by electromagnetic waves}

The total work done per unit time is given, as we have seen, by integrating
$\vec{E} \cdot \vec{j}$ over the whole volume $\tau$:
\begin{equation*}
  \frac{dW}{dt} =  \int_{\tau} d\tau \vec{E} \cdot \vec{j} =
     - \int_{\tau} d\tau \frac{1}{\mu_0} \vec{\nabla} \cdot \Big( \vec{E} \times \vec{B} \Big)
     - \frac{\partial}{\partial t} \int_{\tau} d\tau  \Big( \frac{1}{2\mu_0} |\vec{B}|^2 + \frac{\epsilon_0}{2} |\vec{E}|^2 \Big)
\end{equation*}

\vspace{0.2cm}

Using Gauss' theorem, the volume integral of the divergence of a vector field can be replaced by the
surface integral of that field:
\begin{equation*}
  \frac{dW}{dt} =
     - \oint_{S} d\vec{S} \cdot \frac{1}{\mu_0} \Big( \vec{E} \times \vec{B} \Big)
     - \frac{\partial}{\partial t} \int_{\tau} d\tau  \Big( \frac{1}{2\mu_0} |\vec{B}|^2 + \frac{\epsilon_0}{2} |\vec{E}|^2 \Big)
\end{equation*}

As anticipated, the above tells you that the {\bf work done on the charge q} is, indeed, related with the
{\bf decrease of the energy stored in the e/m field} (2nd term on right-hand side).

\end{frame}


%
%
%
%

\begin{frame}{Energy carried by electromagnetic waves}

But there is another term on the right-hand side of:
\begin{equation*}
  \frac{dW}{dt} =
     - \oint_{S} d\vec{S} \cdot \frac{1}{\mu_0} \Big( \vec{E} \times \vec{B} \Big)
     - \frac{\partial}{\partial t} \int_{\tau} d\tau  \Big( \frac{1}{2\mu_0} |\vec{B}|^2 + \frac{\epsilon_0}{2} |\vec{E}|^2 \Big)
\end{equation*}

It is the surface term:
\begin{equation*}
     \oint_{S} d\vec{S} \cdot \frac{1}{\mu_0} \Big( \vec{E} \times \vec{B} \Big)
\end{equation*}
What does that term correspond to?\\

\vspace{0.2cm}

This term describes the {\em energy carried away from the volume $\tau$,
through its surface S, by the electromagnetic fields}.\\

\vspace{0.2cm}

So the energy stored in the e/m field in volume $\tau$ is decreased, because of the
work done on the charge, but also because energy is flowing out.\\

\vspace{0.2cm}

This is the {\bf Poynting theorem}.

\end{frame}


% ------------------------------------------------------------------------------
% ------------------------------------------------------------------------------
